\section{Elección de ingeniería para tres esquemas de puente}

La clase vuelve a presentar las tres rutas bajo la misma perspectiva sistémica en las diapositivas 41--42. No son enfoques mutuamente excluyentes: el \texttt{augmentation} puede generar direcciones de alta frecuencia, el \texttt{retrieval} preserva experiencias de cola larga y el \texttt{RL} trata estructuras que no pueden recuperarse directamente pero que se pueden reconstruir mediante procesos generales. El verdadero diseño del trabajo consiste en asignar la capacidad de cómputo para entrenamiento, almacenamiento, latencia de inferencia y responsabilidad por recuperación de errores.

\subsection{Marco unificado: convertir \texttt{latent information} en el contexto actual}

Las tres rutas realizan esencialmente la misma tarea: antes de generar una respuesta, hacer que las relaciones necesarias para la tarea ingresen a la computación actual en una forma utilizable. La diferencia radica únicamente en el momento y el soporte. El \texttt{augmentation} explicita y escribe los parámetros antes del entrenamiento; el \texttt{retrieval} recupera desde una memoria externa durante la prueba; el \texttt{RL} genera un conjunto de trabajo interno a partir de los parámetros durante la prueba.

\lecturefigure{slide-041.jpg}{Las tres rutas de puente aprovechan el razonamiento in-context}{Diapositiva oficial Stanford CS25 V6 Lecture 06, p. 41}

\begin{importantbox}{Abstracción unificada: Construcción del Contexto}
Se puede expresar el sistema como $C_q=B(D,M,\theta,q)$: el puente $B$ construye el contexto de trabajo $C_q$ basándose en los datos de entrenamiento $D$, la memoria externa $M$, los parámetros del modelo $\theta$ y la pregunta $q$. Cada uno de los tres métodos coloca la carga computacional principal en la construcción de datos fuera de línea, la recuperación externa en línea o la generación interna en línea.
\end{importantbox}

\subsection{Cómputo, cobertura y dependencias}

Las pros y contras de la diapositiva 42 no deben leerse simplemente como "quién tiene mayor precisión". El \texttt{augmentation} en tiempo de entrenamiento incrementa los costos de generación y entrenamiento, pero abarata el despliegue; el \texttt{retrieval} es gratuito durante el entrenamiento, pero exige un sistema de memoria confiable e incrementa la I/O y los tokens por cada solicitud; el \texttt{RL} puede permitir transferencia entre distribuciones de aumento explícito, pero tanto el entrenamiento como la inferencia requieren cómputo adicional de rollout, siendo débil respecto a la reversibilidad.

\lecturefigure{slide-042.jpg}{Comparación de beneficios y costos del Augmentation, Retrieval y RL}{Diapositiva oficial Stanford CS25 V6 Lecture 06, p. 42}

\begin{knowledgebox}{Lectura de la figura: ninguna columna domina en todas las dimensiones}
Esta página no es un ranking de métodos, sino una frontera de Pareto. El \texttt{augmentation} cambia cómputo de entrenamiento por baja latencia; el \texttt{retrieval} cambia sistema externo por cobertura abierta; el pensamiento \texttt{RL thinking} cambia generación adicional por transferencia de proceso. Al seleccionar, primero se deben fijar las restricciones del negocio y luego buscar soluciones aceptables, en lugar de inferir una arquitectura unificada a partir de la barra más alta de algún benchmark.
\end{knowledgebox}

\begin{table}[H]
\centering
\small
\caption{Tabla de decisiones de ingeniería para los tres esquemas de puente}
\begin{tabularx}{\textwidth}{L{0.18\textwidth}X X X}
\toprule
Dimensión & Train-time augmentation & Test-time retrieval & Test-time RL thinking \\
\midrule
Costo principal & Generación, validación y reentrenamiento fuera de línea & Índice, recuperación y contexto adicional & Rollout de RL y más tokens por generación \\
Latencia de inferencia & Baja, las conclusiones ya están en los parámetros & Medio-alta, depende de la recuperación y el contexto largo & Medio-alta, depende de la longitud del razonamiento \\
Alcance & Limitado por la familia de problemas previamente aumentados & Limitado por la recuperación de memoria y la calidad del episodio & Limitado por el procedimiento aprendido y el espacio de búsqueda \\
Rastreabilidad & Requiere guardar el origen de las muestras aumentadas & Puede citar el episodio original & El CoT interno no siempre es fiel o verificable \\
Fallo típico & alucinación, deriva, omisión de objetivos futuros & Memoria relevante no recuperada o conflictiva & Búsqueda extensa, sobreajuste de recompensa, enumeración de reversal \\
\bottomrule
\end{tabularx}
\end{table}

El costo total se puede expresar aproximadamente como:
$$
C_{\text{total}}=
C_{\text{base-train}}+C_{\text{augment}}+C_{\text{RL}}
+N_{\text{query}}\left(C_{\text{retrieve}}+C_{\text{reason}}\right).
$$
Donde $C_{\text{augment}}$ y $C_{\text{RL}}$ son costos únicos del lado de entrenamiento, mientras que $C_{\text{retrieve}}$ y $C_{\text{reason}}$ se pagan repetidamente según el volumen de consultas $N_{\text{query}}$. Al seleccionar un esquema, es obligatorio calcular junto con el volumen futuro de llamadas, los SLO de latencia y el costo del error.

\subsection{Preguntas y respuestas: contexto largo, escala del modelo y capacidad de recuperación}

En las preguntas y respuestas, Lampinen explica que la estabilidad del ICL depende del modelo y del contenido. Un hecho individual suele ser fácil; incluso en un contexto de cientos de miles de tokens puede tener éxito, pero los entes reales son más fáciles de recuperar que los tokens sin sentido (\texttt{nonsense}), porque el preentrenamiento ya enseña al modelo cómo manejar estos primeros. Por lo tanto, la longitud del contexto y "qué tan familiar está el modelo con las formas de recuperación de este tipo de información" no deben confundirse en una sola variable.

\teachervoice{00:46:14--00:48:05, la advertencia clave que dio el docente es: para datos de la misma longitud, reemplazar nombres reales con palabras sin sentido hace que el rendimiento del contexto disminuya significativamente. La \texttt{tokenización} puede contribuir en parte, pero un factor mayor es la habilidad de recuperación formada durante el preentrenamiento.}

\begin{warningbox}{No prediga la capacidad del ICL basándose en una sola longitud de ventana de contexto}
La ventana permite colocar datos, lo que solo indica que la capacidad de almacenamiento es suficiente; el rendimiento real también está influenciado por la posición, la forma del ente, interferencias de atención, redacción de la consulta, distribución del contexto de duración del entrenamiento y habilidades de recuperación del modelo. Las evaluaciones de despliegue deben utilizar tipos de contenido reales, en lugar de medir solo una plantilla única tipo aguja en un heno.
\end{warningbox}

Sobre los parámetros y la dispersión (\texttt{sparsity}), la clase no ofreció una conclusión clara. Un modelo más grande podría mejorar simultáneamente el rendimiento paramétrico y contextual; la dispersión MoE también podría cambiar la capacidad y las rutas, pero los experimentos presentados no aislaron estas variables. Una conclusión razonable es la necesidad de un estudio de escalado controlado, en lugar de afirmar que la escala eliminaría automáticamente la brecha.

\subsection{Preguntas y respuestas: costos ocultos del augmentation}

El docente estima que el \texttt{offline augmentation} podría ser al menos un orden de magnitud más costoso que un pequeño ajuste fino (\texttt{fine-tuning}) con datos, ya que primero se requiere múltiples rondas de inferencia/generación y luego entrenamiento en un conjunto de datos más grande. Esto es una estimación empírica y no un benchmark de publicación, pero sugiere que los equipos no deben escribir "gratuito en tiempo de inferencia" en sus propuestas e ignorar los costos de GPU fuera de línea, validación e iteración.

\teachervoice{01:08:09--01:09:43, el docente enumera tres tipos de costo: la generación y el entrenamiento hacen que el \texttt{augmentation} sea muy costoso; las muestras generadas pueden alucinar; los datos de ajuste fino más grandes también podrían distorsionar el conocimiento original. La regularización puede mitigar la deriva (\texttt{drift}), pero el compromiso específico depende de la configuración de entrenamiento.}

\begin{knowledgebox}{Plano de control a registrar antes de lanzar Augmentation}
\begin{enumerate}
  \item Modelo de generación, prompt, temperatura y número de muestras;
  \item Documentación fuente y cadena de razonamiento verificable para cada nueva conclusión;
  \item Tasas de filtrado de alucinaciones, duplicados y conflictos;
  \item Comparación antes y después del tarea original, la nueva tarea latente y la capacidad no relacionada;
  \item Tokens adicionales generados, tokens de entrenamiento, FLOPs, tiempo cronometrado y horas de GPU;
  \item Estrategias de KL, decay de pesos, reprodución (\texttt{replay}) o datos híbridos para prevenir deriva de parámetros.
\end{enumerate}
\end{knowledgebox}

\subsection{Preguntas y respuestas: el diseño del prompt es tanto un algoritmo como un riesgo de fuga}

El artículo reutilizó el mismo conjunto de prompts en tareas de relación factual, donde el impacto del ablitación fue menor al esperado; sin embargo, el conferencista aclaró explícitamente que estos prompts vinculan hechos y relaciones de entes, y no necesariamente se trasladan al razonamiento matemático. Si un prompt pide directamente generar la forma objetivo de prueba, esto se acerca a "decirle la respuesta al modelo"; si es demasiado general, podría no ser capaz de desarrollar la relación latente necesaria.

\teachervoice{01:10:03--01:12:08, el docente dice que el prompt pide aproximadamente reescribir el documento y conectar otros entes del corpus, reconociendo al mismo tiempo que es difícil evitar completamente el conocimiento previo de la tarea (\texttt{task foreknowledge}). Los prompts de \texttt{augmentation} universales para múltiples dominios siguen siendo un problema abierto.}

\begin{warningbox}{Problemas de auditoría por fuga en el prompt}
Si se evalúa la reversibilidad, pero se escribe directamente "generar todas las relaciones inversas" en el prompt de \texttt{augmentation}, la mejora del modelo prueba que la construcción de datos dirigida es efectiva, no que se descubrió automáticamente cualquier estructura latente. El reporte debe distinguir entre prompts de conexión genéricos, prompts conscientes de la tarea y fugas de objetivos directos.
\end{warningbox}

\subsection{Resumen de la sección}

Los tres esquemas de puente pueden unificarse como construcción de contexto, pero las posiciones del costo difieren: el \texttt{augmentation} avanza la computación hacia adelante, el \texttt{retrieval} externaliza el almacenamiento y la recuperación, y el \texttt{RL} deja la búsqueda dentro del modelo. La elección para producción debe considerar conjuntamente cobertura, latencia, origen (\texttt{provenance}), recall de recuperación, alucinación, deriva y volumen de consultas, en lugar de observar solo un gráfico de barras de precisión.
